% Source: https://artofproblemsolving.com/community/c3h1708628_geometry_help
% Stable identity: aops-topic-1708628-q01
When $B$ lands on $D$, every point of the crease is equidistant from $B$ and $D$, so the crease $EF$ is the perpendicular bisector of $BD$. Let it meet $BD$ at $O$, the midpoint of $BD$. Since $BD=\sqrt{8^2+6^2}=10$, $OD=5$.

\begin{asy}
usepackage("amsmath");
size(6.5cm);
defaultpen(fontsize(10pt));
pair A=(0,0),B=(8,0),C=(8,6),D=(0,6),O=(4,3),Ep=(1.75,0),F=(6.25,6);
draw(A--B--C--D--cycle,linewidth(0.9));
draw(B--D,blue+dashed);
draw(Ep--F,red+linewidth(1.1));
real q=0.35;pair u=unit(F-O),v=unit(D-O);draw(O+q*u--O+q*(u+v)--O+q*v);
draw((8-q,6)--(8-q,6-q)--(8,6-q));
dot(O);dot(Ep);dot(F);
label("$A$",A,SW);label("$B$",B,SE);label("$C$",C,NE);label("$D$",D,NW);
label("$E$",Ep,S);label("$F$",F,N);label("$O$",O,dir(8));
label("$8$",(4.5,0),S);label("$6$",(8,3),E);
label("\mbox{}",(min(currentpicture,user=true)+max(currentpicture,user=true))/2,black); // keeps the text after the figure black when LaTeX typesets the labels inline
\end{asy}

\textbf{Why triangles $BCD$ and $FOD$ are similar.} Both are right triangles: $\angle BCD=90^\circ$ (corner of the rectangle) and $\angle FOD=90^\circ$ (the crease is perpendicular to $BD$). They also share the angle at $D$, because $O$ lies on $DB$ and $F$ lies on $DC$, so $\angle ODF=\angle BDC$. Two equal angles make the triangles similar, with $O\leftrightarrow C$, $F\leftrightarrow B$, $D\leftrightarrow D$. Hence
\[
\dfrac{OF}{OD}=\dfrac{CB}{CD}=\dfrac{6}{8}\ \Longrightarrow\ OF=\dfrac{3}{4}\cdot5=\dfrac{15}{4}.
\]
By the symmetry of the rectangle about $O$ (a half-turn about $O$ swaps $B,D$ and swaps the lines $AB$, $CD$), $OE=OF$. Therefore
\[
EF=2\cdot\dfrac{15}{4}=\boxed{\dfrac{15}{2}}.
\]
Check with coordinates $A(0,0)$, $B(8,0)$, $D(0,6)$: the crease passes through $E(1.75,0)$ and $F(6.25,6)$, and $EF=\sqrt{4.5^2+6^2}=7.5$.
